%% Allgemeines -----------------------------------

% Fortlaufende Nummerierung - Entweder 1, 2, 3, 4, ... oder mit einem semantischen Schema
\def \decisionID {16}

% Datum, an dem der Record anefertigt wurde
\def \decisionDate {26.11.2026}

% Wer hat an dem Record mitgewirkt?
\def \decisionParticipants {Niklas Käsehagen, OneTutor (Customer Proxy)}

% Der Status in dem sich die Anforderung in diesem Record aktuell befindet.
% Möglichkeiten: Vorgeschlagen, Akzeptiert, Abgelehnt, Abgelöst
\def \decisionStatus {Vorgeschlagen}

% ID des Records, der diesen Record ablöst, falls zutreffend
\def \decisionSuperseded {}

% ID des Records, den dieser Record ablöst, falls zutreffend
\def \decisionSuperseding {}

%% Anforderung -----------------------------------

% Titel der Anforderung
\def \requirementTitle {Der Code, besonders die Ähnlichkeitslogik, sollte gut lesbar und anpassbar sein}

% Kontext der Anforderung
\def \requirementContext {Es gibt Teile in der Software, die nach Abschluss dieses Projekts möglicherweise erneut angepasst werden müssen, weil sie noch nicht fest definiert sind. Ein Beispiel dafür ist die Genauigkeit des Ähnlichkeitsmaßes.}

% Konsequenz der Anforderung
\def \requirementConsequences {Der Code sollte, besonders im Bereich des Ähnlichkeitsmaßes, gut lesbar und ausreichend kommentiert sein, damit er später leicht angepasst werden kann. }
